\documentclass[12pt,a4paper]{article}
\usepackage[utf8]{inputenc}
\usepackage[T2A]{fontenc}
\usepackage[russian]{babel}
\usepackage{amsmath,amssymb,amsthm}
\usepackage{booktabs,array}
\usepackage{geometry}
\usepackage{microtype}
\geometry{margin=2.3cm}

\newtheorem{proposition}{Предложение}
\newtheorem{nogocriterion}{Критерий провала}

\title{Факторные операторы семейного меню:\\
иерархия без выведенного смешивания}
\author{}
\date{5 августа 2026 г.}

\begin{document}
\maketitle

\section{Канонические операторы}

Четыре spin-состояния меню маркируются точками
$F_2^2=\{(p,q)\}$. После удаления тривиального состояния три нетривиальных
характера образуют семейный триплет. Два геометрических фактора носителя
задают независимые трансляции
\begin{equation}
 T_{RP^3}(p,q)=(p+1,q),
 \qquad
 T_{S^1}(p,q)=(p,q+1).
\end{equation}
В базисе
$\{\chi_{RP^3},\chi_{S^1},\chi_{RP^3+S^1}\}$ они диагональны:
\begin{equation}
 T_{RP^3}=\operatorname{diag}(-1,1,-1),
 \qquad
 T_{S^1}=\operatorname{diag}(1,-1,-1).
\end{equation}
Следовательно, в отличие от одного shear-дефекта, факторизация действительно
даёт два канонических коммутирующих оператора и различает все три семейных
направления.

Добавление геометрического shear
$(p,q)\mapsto(p,q+p)$ порождает группу порядка $8$. Её коммутант на триплете
двумерен, то есть полная инвариантность снова оставляет разложение $1+2$.
Поэтому полный порядок не выбирает три различные массы без нарушения части
симметрии.

\section{Проверка естественных ядер}

Общий факторный лапласиан имеет вид
\begin{equation}
 L(w_3,w_1)=w_3(I-T_{RP^3})+w_1(I-T_{S^1})
\end{equation}
и спектр
\begin{equation}
 \{2w_3,2w_1,2(w_3+w_1)\}.
\end{equation}
Топология не фиксирует функцию, переводящую длину цикла в Yukawa-вес.
Проверены три простейших выбора: обратная длина, обратный квадрат и
туннельный вес. Ни один не воспроизводит наблюдаемую трёхуровневую иерархию.

Наиболее сильным из проверенных оказался не лапласиан, а аддитивное winding-
правило
\begin{equation}
 m_{p,q}\propto \exp(-\pi p-2\pi q),
\end{equation}
которое после нормировки даёт
\begin{equation}
 (e^{-2\pi},e^{-\pi},1)
 =(0.00186744,0.0432139,1).
\end{equation}
Для down-кварков его логарифмический RMS равен $0.591$, для заряженных
лептонов $1.342$, а для up-кварков $3.754$. Поэтому найден общий масштабный
рисунок, но не универсальная массовая формула.

\section{Mixing gate}

Любые два Yukawa-оператора вида
\begin{equation}
 Y_u=f_u(T_{RP^3},T_{S^1}),
 \qquad
 Y_d=f_d(T_{RP^3},T_{S^1})
\end{equation}
диагонализуются в одном и том же характерном базисе. Отсюда следует
\begin{equation}
 V_{CKM}=I.
\end{equation}
SU(5)-чётность может менять собственные значения, но не общий левый семейный
базис.

\begin{nogocriterion}[Провал чисто факторной семейной алгебры]
Коммутирующие трансляции решают задачу различимости трёх семейных направлений,
но не выводят ни kernel длина--Yukawa, ни ненулевое смешивание. Следующая
допустимая конструкция должна получить из действия секторно-зависимую
некоммутирующую shear-вставку; вводить её коэффициент по CKM-данным запрещено.
\end{nogocriterion}

Полный машинный протокол сохранён в
\texttt{s2t\_family\_factor\_operator\_audit.py} и
\texttt{s2t\_family\_factor\_operator\_results.json}.

\end{document}